\section{What Is Copied and What Is Shared}
\label{sec:ch12-copied-and-shared}
\index{duplication!the line worth drawing|(}%

Chapter~\ref{ch:the-first-cut-extracting-catalog} duplicated one file on purpose
and said the duplication was the point. That is an easy position to hold once. At
four copies a reader is entitled to ask whether it is still a principle or just a
habit, so here is the whole ledger at tag \code{ch12}:

\begin{minted}[fontsize=\footnotesize]{text}
ProductKey.TryDecode      3 copies   pricing, inventory, reviews
CustomerKey.TryDecode     1 copy     accounts; ProductKey, one name changed
Product stub              4 copies   three that resolve one, one that cannot
Customer stub             2 copies   beside the real Customer in accounts
ProductCategory           1 copy     pricing, beside Catalog's
Money                     2 copies   pricing and ordering, both @shareable
PageCursorExtension       2 copies   the two subgraphs with a connection
\end{minted}

Fifteen files that exist because another file exists. The test I applied to
every one of them is the same question: \textbf{do two services have to agree
about this?} If they do, it is a contract, and a contract gets copied. If they
do not, it is platform, and platform gets shared.

\index{key format!an agreement rather than a class}%
\code{ProductKey.TryDecode} is the clearest case of the first kind. It encodes
one subgraph's belief about how another subgraph spells a product. A shared
class would make that belief a compile-time fact, and the entire value of a
federated graph is that it is not one: the moment those five teams cannot deploy
without agreeing on a build, the extraction has bought nothing.

That argument has a cost and it is worth naming rather than skating past.
Five files parse against \code{typeof(Guid)} and check a type name: the four
that decode a product key, and Accounts decoding a customer one. If Mosaic ever
changed how it encodes identifiers, five files in five repositories owned by
five teams would have to change, and \emph{there is no build that would fail if
one of them did not}. The decoupling is paid for in coordination. What
catches it instead is the gate, which now asks Pricing to price every key
Catalog handed out and asks Accounts to resolve every key Reviews handed out.
Those two assertions exist for exactly this and nothing else.

\subsection*{The copy the composer actually inspects}

\index{shareable@\code{@shareable}!on a value type}%
\code{Money} is the third kind of duplication and the only one the composer
looks at field by field. A key format is a string both sides agree to parse. An
enum is a list of names. \code{Money} is an object type that two subgraphs
declare, and Apollo Federation refuses that outright unless somebody says the
fields may be resolved in more than one place:

\begin{minted}[fontsize=\footnotesize]{csharp}
[Shareable]
public sealed record Money(decimal Amount, string Currency);
\end{minted}

The attribute goes on the class, and \code{ShareableAttribute.TryConfigure}
applies it to the object type descriptor, which marks every field. Marking every
field is what matters, because the composer checks one field at a time. Strip
the attribute off both records and the complaint comes back per field rather
than per type: a heading naming \code{Money}, then two lines for \code{amount}
and two more for \code{currency}, all four naming \code{pricing} and
\code{ordering}. You forgot one attribute on one record and the composer hands
you back four complaints about fields.

\index{shareable@\code{@shareable}!what it does not promise}%
It is worth being precise about what \code{@shareable} is not, because the name
oversells it. It does not merge the two declarations and it does not make the
router prefer either. It says the graph may take this field from whichever
subgraph it is already talking to. Which is fine here and would not be fine
everywhere, and the reason is visible in what the two services mean by the type:

\begin{minted}[fontsize=\footnotesize,breaklines]{text}
Pricing's Money    what a product costs now
Ordering's Money   what a customer paid on the day
\end{minted}

Two services sharing a value type do not necessarily share anything else, and
nothing in composition can tell which situation it is looking at. This is
chapter~\ref{ch:entity-resolution-done-right}'s argument about \code{@provides}
arriving in a different shape: a value two services can both produce is a value
two services can disagree about.

\subsection*{The library, and what is deliberately not in it}

\index{Mosaic.ServiceDefaults@\code{Mosaic.ServiceDefaults}}%
\code{src/Mosaic.ServiceDefaults} holds the other side of the test. The request
timeline, the diagnostic listener, the pipeline report, the two counters, the
snake\_case convention, and four builder extensions, of which
section~\ref{sec:ch12-what-a-service-costs-to-make} showed one. None of those is
anything two services have to agree about; they are things all six should do the
same way so that six logs can be read by one person.

One line in its \code{csproj} is worth a sentence because it will catch
somebody:

\begin{minted}[fontsize=\footnotesize]{text}
<PackageReference Include="Microsoft.EntityFrameworkCore.Relational" />
\end{minted}

\code{DbCommandInterceptor} is a relational type, not a core one. The six
services never name that package because the Npgsql provider brings it in; a
library with no provider has to ask for it, and the error until it does is
\code{CS0246} on a type that looks like it should be in the package the project
already references.

\index{team contracts!the two lists}%
If I were handing this to five teams tomorrow, the two lists are the whole
handover document. \emph{Shared, and you take a version bump when it changes}:
the platform library, the compose file, the router config, the key format's
shape. \emph{Yours, and nobody else's business}: your database, your seed data,
your DataLoaders, your resolvers, and which fields you contribute to somebody
else's entity. The second list is longer than the first and that is the test of
whether a split was worth doing.

\medskip
Which leaves the part of a migration that cannot be done by moving files: taking
over a field that is already answering traffic.

\index{duplication!the line worth drawing|)}%
